%% RQ5 results (Phase 8). Every number is a macro from tables/numbers.tex.
We re-examined five published comparisons of releases (Section~\ref{sec:claims}): Rosay et al.\ on LYCOS-IDS2017 against the original \cicids{} (R1) \cite{rosay2021}, Engelen et al.\ on WTMC-2021 (R2) \cite{engelen2021}, Lanvin et al.\ on the CRiSIS-2022 variants (R3) \cite{lanvin2023}, Sarhan et al.\ on the NetFlow version of UNSW-NB15 (R4) \cite{sarhan2022} and Pekar and Jozsa on five releases per day (R5) \cite{pekar2024}. They give \RQfiveRows{} headline contrasts, listed in Table~\ref{tab:rq5} and drawn in Fig.~\ref{fig:rq5}. Each contrast is examined in three steps: the per-flow comparison rebuilt the way the authors made it, a decomposition of that per-flow gap, and the same contrast on conversations present in both releases, with one ground truth, the conversation-grouped split and conversation scoring.

% Generated by 03_Code/make_numbers.py from 04_Results/phase8/verdicts_sched.csv. Do not edit by hand.
\begin{sidewaystable}
\caption{Published claims that one release is better than another. Contrast: release B vs release
A; every gap is B minus A in the claim's metric (F1, recall and MCC without a class: attack against
benign). R5: (a) Tuesday, each refined release against CIC original; (b) Thursday; (c) the claim
of no difference, NFS-2023-TE against NFS-2023-nTE on each day. Published: the gap the authors
report. Per flow, as published: our rebuild of their comparison. Per flow, shared conv.: the
rebuilt per-flow comparison restricted to the conversations present in both releases, with schedule
labels and the grouped split. Per conversation: the same comparison scored once per conversation;
the verdict is read from this column. Attempts benign: the verdict when attempted attacks are folded
into benign (Section~\ref{sec:gt}); n/a, UNSW-NB15 has no attempted attacks; not computed, the
contrast involves no class with attempted attacks, or was not rerun. R4: the per-flow rebuild uses
NF-UNSW-NB15-v2 itself and, on the second line, NF-UNSW-NB15-v3; the shared-conversation and
conversation-level contrasts use NF-UNSW-NB15-v3, which has the timestamps that v2 lacks.
Prevalence: the published and per-flow gaps use each release's test data as the authors split it;
the shared-conversation and per-conversation gaps use the subsample of Section~\ref{sec:protocol},
in which attacks are more prevalent. Recall (R3, and R4 Fuzzers and Worms) does not depend on class
prevalence; F1 and MCC do. R1 models: RF random forest, DT decision tree, LDA linear
discriminant analysis. Brackets: 95 percent paired cluster bootstrap intervals over conversations.
Verdicts: zero by construction, the two releases are one arm under schedule labels; data not
obtained, the release could not be downloaded. R5 rows name the capture day after the
metric.}\label{tab:rq5}
\footnotesize
\setlength{\tabcolsep}{3pt}
\begin{tabularx}{\linewidth}{@{}Llr*{3}{>{\raggedright\arraybackslash}p{2.4cm}}*{2}{>{\raggedright\arraybackslash}p{1.6cm}}@{}}
\toprule
Contrast & Metric & Published & Per flow, as published & Per flow, shared conv. & Per conversation & Verdict & Attempts benign \\
\midrule
\multicolumn{8}{@{}l}{\emph{R1, Rosay et al.}} \\
LYCOS vs CIC original, RF & MCC & 0.001 & 0.001 [0.001, 0.001] & 0.002 [0.001, 0.002] & 0.001 [0.000, 0.001] & survives & not computed \\
LYCOS vs CIC original, DT & MCC & 0.001 & 0.001 [0.000, 0.001] & 0.001 [0.001, 0.001] & 0.001 [0.001, 0.001] & survives & not computed \\
LYCOS vs CIC original, LDA & MCC & 0.155 & 0.155 [0.152, 0.158] & 0.210 [0.207, 0.213] & 0.164 [0.162, 0.167] & survives & not computed \\
\addlinespace
\multicolumn{8}{@{}l}{\emph{R2, Engelen et al.}} \\
WTMC-2021 vs CIC original & F1, Bot & 0.390 & 0.183 [0.168, 0.199] & 0.048 [0.039, 0.057] & 0.030 [0.021, 0.040] & survives & survives \\
WTMC-2021 vs CIC original & F1, Web brute force & 0.180 & 0.222 [0.186, 0.251] & 0.013 [0.000, 0.026] & $-0.040$ [$-0.055$, $-0.027$] & reverses & not supported \\
WTMC-2021 vs CIC original & F1, Web XSS & 0.610 & 0.458 [0.254, 0.586] & 0.040 [0.004, 0.072] & 0.197 [0.153, 0.239] & survives & not estimable \\
WTMC-2021 vs CIC original & F1, Infiltration & 0.100 & 0.190 [0.063, 0.424] & n/a & n/a & not estimable & not estimable \\
WTMC-2021 vs CIC original & F1, Heartbleed & 0.160 & 0.040 [0.040, 0.040] & n/a & n/a & not estimable & not estimable \\
WTMC-2021 vs CIC original & F1, Web SQL injection & $-0.110$ & 0.214 [$-0.138$, 0.500] & 0.043 [$-0.237$, 0.241] & $-0.234$ [$-0.505$, $-0.027$] & not estimable & not estimable \\
\addlinespace
\multicolumn{8}{@{}l}{\emph{R3, Lanvin et al.}} \\
CRiSIS-2022 Exp\_RP vs CRiSIS-2022 Exp\_RPD & Recall & 0.200 & 0.211 [0.209, 0.212] & 0.000 [0.000, 0.000] & 0.000 [0.000, 0.000] & not supported & survives \\
CRiSIS-2022 Exp\_RP vs CRiSIS-2022 Exp\_RPD & Recall & 0.000 & $-0.006$ [$-0.006$, $-0.006$] & 0.000 [0.000, 0.000] & 0.000 [0.000, 0.000] & survives & survives \\
CRiSIS-2022 Exp\_RPD vs CRiSIS-2022 Exp\_RD & Recall & $-0.200$ & $-0.217$ [$-0.218$, $-0.215$] & 0.000 [0.000, 0.000] & 0.000 [0.000, 0.000] & zero by construction & zero by construction \\
\bottomrule
\end{tabularx}
\end{sidewaystable}

\begin{sidewaystable}
\addtocounter{table}{-1}
\renewcommand{\theHtable}{\arabic{table}.cont}
\caption{Published claims that one release is better than another (continued).}
\footnotesize
\setlength{\tabcolsep}{3pt}
\begin{tabularx}{\linewidth}{@{}Llr*{3}{>{\raggedright\arraybackslash}p{2.4cm}}*{2}{>{\raggedright\arraybackslash}p{1.6cm}}@{}}
\toprule
Contrast & Metric & Published & Per flow, as published & Per flow, shared conv. & Per conversation & Verdict & Attempts benign \\
\midrule
\multicolumn{8}{@{}l}{\emph{R4, Sarhan et al.}} \\
NF-UNSW-NB15-v3 vs UNSW-NB15 original & F1 & 0.050 & v2 $-0.009$ [$-0.016$, 0.004]\newline v3 0.022 [0.016, 0.036] & $-0.019$ [$-0.031$, 0.013] & $-0.003$ [$-0.004$, $-0.002$] & reverses & n/a \\
NF-UNSW-NB15-v3 vs UNSW-NB15 original & Recall, Fuzzers & 0.301 & v2 0.254 [0.214, 0.296]\newline v3 0.267 [0.226, 0.311] & $-0.022$ [$-0.027$, $-0.016$] & $-0.019$ [$-0.024$, $-0.013$] & reverses & n/a \\
NF-UNSW-NB15-v3 vs UNSW-NB15 original & Recall, Worms & 0.807 & v2 0.672 [0.613, 0.734]\newline v3 0.608 [0.520, 0.692] & 0.685 [0.594, 0.763] & 0.685 [0.593, 0.763] & survives & n/a \\
\addlinespace
\multicolumn{8}{@{}l}{\emph{R5, Pekar and Jozsa}} \\
(a) WTMC-2021 vs CIC original & F1, Tuesday & 0.129 & 0.129 [0.122, 0.136] & 0.105 [0.101, 0.108] & 0.106 [0.100, 0.111] & survives & not computed \\
(a) CRiSIS-2022 Exp\_RP vs CIC original & F1, Tuesday & 0.128 & 0.129 [0.122, 0.136] & 0.105 [0.101, 0.109] & 0.105 [0.100, 0.111] & survives & not computed \\
(a) NFS-2023-nTE vs CIC original & F1, Tuesday & 0.128 & 0.128 [0.122, 0.136] & 0.105 [0.101, 0.109] & 0.104 [0.098, 0.110] & survives & not computed \\
(a) NFS-2023-TE vs CIC original & F1, Tuesday & 0.128 & 0.129 [0.122, 0.136] & $-0.013$ [$-0.017$, $-0.009$] & 0.104 [0.098, 0.110] & survives & not computed \\
(b) NFS-2023-nTE vs NFS-2023-TE & F1, Thursday & 0.050 & 0.050 [0.047, 0.052] & 0.177 [0.168, 0.186] & 0.292 [0.279, 0.305] & survives & not computed \\
(b) CRiSIS-2022 Exp\_RP vs WTMC-2021 & F1, Thursday & 0.047 & 0.047 [0.008, 0.099] & 0.000 [$-0.001$, 0.001] & 0.000 [$-0.001$, 0.001] & not supported & not computed \\
(b) NFS-2023-nTE vs CIC original & F1, Thursday & 0.051 & 0.050 [0.037, 0.064] & 0.033 [0.028, 0.037] & 0.015 [0.010, 0.020] & survives & not computed \\
(c) NFS-2023-TE vs NFS-2023-nTE & F1, Tuesday & 0.000 & 0.000 [$-0.001$, 0.001] & $-0.118$ [$-0.120$, $-0.116$] & 0.000 [$-0.001$, 0.000] & survives & not computed \\
(c) NFS-2023-TE vs NFS-2023-nTE & F1, Wednesday & 0.000 & 0.000 [0.000, 0.000] & $-0.044$ [$-0.045$, $-0.042$] & $-0.022$ [$-0.023$, $-0.021$] & contradicted & not computed \\
(c) NFS-2023-TE vs NFS-2023-nTE & F1, Thursday & $-0.050$ & $-0.050$ [$-0.052$, $-0.047$] & $-0.177$ [$-0.186$, $-0.168$] & $-0.292$ [$-0.305$, $-0.279$] & contradicted & not computed \\
(c) NFS-2023-TE vs NFS-2023-nTE & F1, Friday & 0.000 & 0.000 [0.000, 0.000] & $-0.006$ [$-0.006$, $-0.006$] & 0.000 [0.000, 0.000] & survives & not computed \\
\bottomrule
\end{tabularx}
\end{sidewaystable}


\paragraph{The per-flow rebuilds reproduce most published gaps.} Where we could follow the authors' code, the rebuild lands on their numbers: for R1 and R5 every rebuilt score is within \RQfiveFidelityMax{} of the published one. R3 reproduces the published fall in recall when the corrected labels are used (\RQfiveQtwoFlow{} against about \RQfiveQtwoPub{}). Two rebuilds differ on one side. In R2 our scores on WTMC-2021 are close to the published ones, but on the original release they are higher (Bot F1 \RQfiveRtwoBotOrigOurs{} against \RQfiveRtwoBotOrigPub{}), so the rebuilt gap is smaller (Bot \RQfiveBotFlow{} against \RQfiveBotPub{}); we did not read the authors' hyperparameters or feature list, which is the likely source. In R4 our binary F1 on the original UNSW-NB15 is \RQfiveRfourOrigOurs{} against a published \RQfiveRfourOrigPub{}, while the NetFlow side matches, so the published binary advantage is not reproduced (gap \RQfiveBinFlow{}, interval \RQfiveBinFlowLo{} to \RQfiveBinFlowHi{}); the per-class gain for Fuzzers is reproduced (\RQfiveFuzzFlow{} against \RQfiveFuzzPub{}). Every deviation from the authors' setups is ours and is listed with the code. Where a rebuild does not reproduce the published baseline, a changed verdict could come from the rebuild rather than from the scoring unit. Table~\ref{tab:rq5} therefore gives, beside the conversation-level gap, the rebuilt per-flow gap restricted to the shared conversations with one truth. Comparing it with the per-flow gap as published separates what changes with the population and the labels from what changes with the unit, within the same rebuild.

\paragraph{Class weights never explain a whole gap.} A per-flow gap between two releases splits into a class weight part, from the change in how many test records each class has, a within-class part, and a part from classes that only one release has. The class weight part is small and mostly negative in R2 (Bot \RQfiveBotC{}, Web XSS \RQfiveXssC{}, Web brute force \RQfiveWebBfC{}) and is zero by design in R1, whose test sets hold equal numbers of attack records per class. In R3 it is \RQfiveQtwoC{} of the \RQfiveQtwoFlow{} gap, the rest within class (\RQfiveQtwoW{}); the flows of the two releases are the same, and what changes is which test records count as attacks. In R4 the two parts have opposite signs (class weights \RQfiveBinC{}, within class \RQfiveBinW{}). The largest structural effect is in R5 on Thursday, where NFS-2023-nTE beats CIC original by \RQfiveThuNteCicPub{} as published. The classes both releases share contribute \RQfiveThuNteCicShared{} to that gap (class weights \RQfiveThuNteCicC{}, within class \RQfiveThuNteCicW{}), and the port scan that NFS-2023-nTE labels on Thursday, a class the original release does not have on that day, contributes \RQfiveThuNteCicNonshared{}. The published advantage comes from labelling an easy class, not from better detection of the classes both releases share.

\paragraph{Verdicts per conversation.} Of the \RQfiveRows{} contrasts, \RevFiveUndecided{} cannot be decided on the same conversations with one truth: \RQfiveNotEstimable{} are not estimable because a class has fewer than \SetMinTest{} test conversations, and \RQfiveZero{} is zero by construction. Of the \RevFiveDecided{} that can be decided, \RQfiveSurvive{} survive (\RevFiveSurvivePct{} percent), \RQfiveReverse{} reverse, \RQfiveContradicted{} claims of no difference are contradicted and \RQfiveNotSupported{} are not supported. The verdict rules were fixed before any conversation-level number was computed: a claim survives when the gap has the published sign and its interval excludes zero, reverses when the interval lies entirely on the other side, and is not supported when the interval includes zero; a claim of no difference survives when the interval lies within \setnum{0.01} of zero and is contradicted when it excludes zero.

Where a claim survives, its size often shrinks. In R1 the LYCOS advantage for linear discriminant analysis holds (\RQfiveLdaConv{}, interval \RQfiveLdaConvLo{} to \RQfiveLdaConvHi{}, against \RQfiveLdaPub{} published). In R2 the WTMC-2021 gain for Bot falls from \RQfiveBotPub{} published to \RQfiveBotConv{} (\RQfiveBotConvLo{} to \RQfiveBotConvHi{}), and for Web XSS from \RQfiveXssPub{} to \RQfiveXssConv{}. The R2 gain for Web brute force reverses (\RQfiveWebBfConv{}, \RQfiveWebBfConvLo{} to \RQfiveWebBfConvHi{}), but this verdict depends on how attempted attacks are labelled: with attempts folded into benign the gap is \RQfiveFoldWebBfConv{} (\RQfiveFoldWebBfConvLo{} to \RQfiveFoldWebBfConvHi{}) and the claim is \RQfiveFoldWebBfVerdict{}, while Bot still \RQfiveFoldBotVerdict{} and Web XSS becomes \RQfiveFoldXssVerdict{}. Infiltration, Heartbleed and Web SQL injection are not estimable in R2.

R3 compares releases that differ only in labels. Under one schedule truth the corrected release and the release it corrects become the same arm, so their gap is zero by construction; the claim cannot be tested this way, and the per-flow result above is the finding. The two R3 contrasts that involve Exp\_RP, the variant without duplicated packets, show the same thing from the other side. Per flow, Exp\_RP raises recall over Exp\_RPD by \RQfiveQthreeFlow{} (\RQfiveQthreeFlowLo{} to \RQfiveQthreeFlowHi{}), close to the published \RQfiveQthreePub{}. Per conversation the gain is \RQfiveQthreeConv{} (\RQfiveQthreeConvLo{} to \RQfiveQthreeConvHi{}): both releases detect about \RQfiveQthreeRecall{} percent of attack conversations, so the published gain comes from how records are counted and labelled, not from what is detected, and the claim is not supported. The claim that Exp\_RD and Exp\_RP give the same recall survives (\RQfiveQfourFlow{} per flow, \RQfiveQfourConv{} per conversation).

R4 changes most. NF-UNSW-NB15-v2 has no timestamps, so it cannot be put on our clock, and the conversation-level contrast uses NF-UNSW-NB15-v3, the same feature set with millisecond timestamps, in its place. On the same conversations, the NetFlow arm is not better at binary detection (\RQfiveBinConv{}, \RQfiveBinConvLo{} to \RQfiveBinConvHi{}, against \RQfiveBinPub{} published) and is worse on Fuzzers (\RQfiveFuzzConv{}, \RQfiveFuzzConvLo{} to \RQfiveFuzzConvHi{}, against \RQfiveFuzzPub{}); only Worms keeps a large gain (\RQfiveWormsConv{}, \RQfiveWormsConvLo{} to \RQfiveWormsConvHi{}, against \RQfiveWormsPub{}). The proxy matters for the binary row. Per flow, NF-UNSW-NB15-v3 beats the original by \RQfiveBinVthreeFlow{} (\RQfiveBinVthreeFlowLo{} to \RQfiveBinVthreeFlowHi{}), in the published direction, so the reversal on this row is a change between the per-flow and the per-conversation comparison on NF-UNSW-NB15-v3, not a property of NF-UNSW-NB15-v2. Because our per-flow rebuild of the original release scores higher than published, the R4 reversals rest on our rebuild and on a proxy release, and we give them less weight than the R1, R3 and R5 verdicts, whose baselines reproduce.

In R5 the Tuesday gains of the refined releases over the original survive at a smaller size (WTMC-2021 \RQfiveTueWtmcConv{} against \RQfiveTueWtmcPub{} published, and CRiSIS-2022, which Pekar and Jozsa took from Exp\_RP, \RQfiveTueCrisisConv{} against \RQfiveTueCrisisPub{}). On Thursday the CRiSIS-2022 advantage over WTMC-2021 is rebuilt per flow (\RQfiveThuCrisisFlow{}, \RQfiveThuCrisisFlowLo{} to \RQfiveThuCrisisFlowHi{}, against \RQfiveThuCrisisPub{} published) but vanishes per conversation (\RQfiveThuCrisisConv{}, \RQfiveThuCrisisConvLo{} to \RQfiveThuCrisisConvHi{}), so it is not supported. The Thursday gap between the two NFStream releases grows each time the comparison becomes defined: \RQfiveThuNteTeFlow{} per flow as published, \RQfiveThuNteTeFlowDef{} (\RQfiveThuNteTeFlowDefLo{} to \RQfiveThuNteTeFlowDefHi{}) per flow on the shared conversations with one truth, and \RQfiveThuNteTeConv{} (\RQfiveThuNteTeConvLo{} to \RQfiveThuNteTeConvHi{}) per conversation. The advantage of NFS-2023-nTE over the original on Thursday survives but shrinks from \RQfiveThuNteCicPub{} to \RQfiveThuNteCicConv{}. Pekar and Jozsa also concluded that TCP flag based expiry does not adversely affect detection. Per conversation this holds on Tuesday and Friday and is contradicted on Wednesday (\RQfiveWedTeConv{}, \RQfiveWedTeConvLo{} to \RQfiveWedTeConvHi{}), where the per-flow gap is zero, and on Thursday (\RQfiveThuTeConv{}, \RQfiveThuTeConvLo{} to \RQfiveThuTeConvHi{}). The Thursday gap comes from false alarms, not missed attacks, and the false-alarm rate, unlike precision and F1, does not depend on the class prevalence of the subsample: recall is \RQfiveThuRecallTe{} for TE and \RQfiveThuRecallNte{} for nTE, but TE flags \RQfiveThuFprTe{} percent of benign test conversations against \RQfiveThuFprNte{} percent for nTE. This has the pattern of the RQ4 residual, benign conversations made false alarms by one flagged record; we checked which records are flagged for the RQ4 classes, not for this contrast.

\paragraph{Answer to RQ5.} Of the \RevFiveDecided{} published cross-release contrasts that can be decided, \RQfiveSurvive{} survive a comparison in which both releases are scored on the same conversations with one truth, and where they survive they are often much smaller. Three reverse, among them the NetFlow version's binary advantage on UNSW-NB15, which rests on a proxy release and on a rebuild that does not reproduce the published baseline, and the Web brute force gain of WTMC-2021, which is no longer a reversal when attempted attacks are labelled benign. The claim that flag-based expiry does not hurt detection fails on two of four days. Two more, the gain from dropping duplicated packets in CRiSIS-2022 and that release's Thursday advantage over WTMC-2021, hold per flow and vanish per conversation. In no claim does the change in class weights explain the whole per-flow gap; the largest single source we found is a class that one release labels and the other does not.
